%
% File: app0A.tex
% Author: Tengxiang Li
%
\chapter{Appendix A: Source Code Snippets}

\section*{A.1 Go Backend: Generative AI Prompt Builder}
The following Go snippet illustrates how the backend constructs system instructions and context boundaries before dispatching the request to the Google Gemini API.

\begin{verbatim}
package ai

import "fmt"

func BuildSummaryPrompt(text string, format string) string {
    systemPrompt := `You are an expert academic tutor. 
    Analyze the provided content and summarize it. 
    Format the output strictly according to the user's format.`
    
    constraint := fmt.Sprintf("Required Format: %s. Do not add conversational filler.", format)
    
    return fmt.Sprintf("%s\n%s\n\n[CONTENT]\n%s", systemPrompt, constraint, text)
}
\end{verbatim}

\section*{A.2 React Frontend: WebSocket Connection Hook}
This custom React hook demonstrates how the frontend establishes a persistent WebSocket connection to receive real-time updates regarding AI generation job statuses.

\begin{verbatim}
import { useEffect, useState } from 'react';

export function useJobSocket(jobId: string) {
    const [status, setStatus] = useState('pending');

    useEffect(() => {
        const ws = new WebSocket(`ws://localhost:8080/ws/jobs/${jobId}`);
        
        ws.onmessage = (event) => {
            const data = JSON.parse(event.data);
            setStatus(data.status);
            if (data.status === 'completed') {
                ws.close();
            }
        };

        return () => ws.close();
    }, [jobId]);

    return status;
}
\end{verbatim}



\chapter{Appendix B: User Manual}

\section*{B.1 Introduction}
This appendix provides a brief manual for end-users operating the Lectura platform. The user interface was designed with standard Material/Tailwind paradigms to minimize the learning curve.

\section*{B.2 Getting Started}
\begin{enumerate}
    \item \textit{Registration / Login}: Navigate to the homepage and click "Sign In". You may authenticate via traditional Email/Password credentials or by using the single sign-on Google OAuth provider.
    \item \textit{Dashboard Overview}: Upon authenticating, you are redirected to the \textit{My Library} dashboard. Here, statistics regarding total study time and recent flashcard activity are displayed.
\end{enumerate}

\section*{B.3 Creating Study Materials}
\begin{enumerate}
    \item Click the \textbf{"New Material"} button in the sidebar.
    \item \textit{Input Source}: Choose between pasting a YouTube video URL, uploading a document (PDF, DOCX), or pasting raw text.
    \item \textit{Configuration}: Select output preferences (e.g., Cornell Notes formatting, Quiz difficulty level).
    \item \textit{Generation}: Click "Generate". A progress indicator will display real-time updates as the backend worker pool processes the request.
\end{enumerate}

\section*{B.4 Interacting with Content}
Once the generation is complete, you will be redirected to the Content Viewer suite:
\begin{itemize}
    \item \textit{Summary Tab}: View and read the generated notes.
    \item \textit{Quiz Tab}: Take a multiple-choice test based on the material. Feedback is provided instantly upon submission.
    \item \textit{Flashcards Tab}: Enter a focused study session to review vocabulary and core concepts using spaced repetition UI patterns.
    \item \textit{Tutor Chat}: Found at the bottom right of the screen, open this panel to ask follow-up questions about the active document. The AI acts as a tutor grounding its answers strictly in the provided text.
\end{itemize}

\section*{B.5 Exporting Data}
To save materials for offline revision, click the \textbf{"Export PDF"} button on the Summary tab. The system will compile the notes and quiz questions into a cleanly formatted academic PDF and trigger a browser download.
